% ============================================================
%   Appendix B  --  Observation Space Breakdown
% ============================================================

\chapter{Observation Space Breakdown}\label{app:observation}

This appendix describes every one of the \(739\) dimensions of the observation vector
\(s_t\). Indices are zero-based and match the slicing conventions used in
\texttt{src/CARLA/Env/carla\_env.py}.

\section{LIDAR block (dimensions 0\textendash 719)}\label{app:obs-lidar}

The first \(720\) dimensions contain three semantic LIDAR channels of \(240\) rays each. Each ray
is normalised to \([0, 1]\), where \(1\) means a clear ray over the whole \(50\,\mathrm{m}\)
range and values close to \(0\) mean an obstacle within \({\approx}\,1\,\mathrm{m}\). Ray
\(k\in\{0,\ldots,239\}\) points towards the azimuth angle \(\phi_k = k \cdot 1.5^\circ\)
measured counter-clockwise from the vehicle's forward direction.

\begin{table}[H]
  \centering
  \caption{LIDAR block of the observation vector.}
  \label{tab:obs-lidar}
  \small
  \begin{tabular}{@{}lll@{}}
    \toprule
    \textbf{Indices} & \textbf{Channel} & \textbf{Description} \\
    \midrule
    \([0, 240)\)   & Combined  & All obstacles (vehicles, pedestrians, static, road edges). \\
    \([240, 480)\) & Dynamic   & Only vehicles and pedestrians. \\
    \([480, 720)\) & Static    & Only walls, barriers, poles and similar static objects. \\
    \bottomrule
  \end{tabular}
\end{table}

\section{Lane features (dimensions 720\textendash 727)}\label{app:obs-lane}

The next \(8\) dimensions are lane features computed from CARLA's Waypoint API at the
vehicle's current position.

\begin{table}[H]
  \centering
  \caption{Lane features.}
  \label{tab:obs-lane}
  \small
  \begin{tabular}{@{}cllp{6cm}@{}}
    \toprule
    \textbf{Idx.} & \textbf{Name} & \textbf{Range} & \textbf{Meaning} \\
    \midrule
    720 & \texttt{lateral\_offset\_norm}    & \([-1, 1]\) & Signed lateral offset of the
        vehicle with respect to the lane centre, normalised by the lane half-width. \\
    721 & \texttt{heading\_error\_norm}     & \([-1, 1]\) & Signed heading error between the
        vehicle and the lane tangent, normalised by \(180^\circ\). \\
    722 & \texttt{on\_edge\_warning}        & \([0, 1]\) & Continuous warning that grows
        quadratically as the closest lane edge approaches. \\
    723 & \texttt{lane\_width\_norm}        & \([0, 1]\) & Lane width at the current
        position, normalised by a reference width of \(4.5\,\mathrm{m}\). \\
    724 & \texttt{dist\_left\_edge\_norm}   & \([0, 1]\) & Distance to the left lane edge,
        normalised by the lane width. \\
    725 & \texttt{dist\_right\_edge\_norm}  & \([0, 1]\) & Distance to the right lane edge,
        normalised by the lane width. \\
    726 & \texttt{lane\_change\_left}        & \(\{0, 1\}\) & \(1\) if a left-side lane change is
        permitted at the current position (only the left direction is exposed). \\
    727 & \texttt{road\_curvature\_norm}    & \([-1, 1]\) & Estimated signed road curvature,
        normalised by a reference curvature. \\
    \bottomrule
  \end{tabular}
\end{table}

\section{Lane-marking type (dimensions 728\textendash 731)}\label{app:obs-marking}

Four binary flags classify the lane-marking type on each side of the vehicle, derived from
CARLA's \texttt{LaneMarkingType} at the current position.

\begin{table}[H]
  \centering
  \caption{Lane-marking type features.}
  \label{tab:obs-marking}
  \small
  \begin{tabular}{@{}cllp{6cm}@{}}
    \toprule
    \textbf{Idx.} & \textbf{Name} & \textbf{Range} & \textbf{Meaning} \\
    \midrule
    728 & \texttt{solid\_left}   & \(\{0, 1\}\) & \(1\) if the left lane marking is solid
        (crossing is forbidden). \\
    729 & \texttt{solid\_right}  & \(\{0, 1\}\) & \(1\) if the right lane marking is solid. \\
    730 & \texttt{dashed\_left}  & \(\{0, 1\}\) & \(1\) if the left marking is dashed
        (lane change permitted). \\
    731 & \texttt{dashed\_right} & \(\{0, 1\}\) & \(1\) if the right marking is dashed. \\
    \bottomrule
  \end{tabular}
\end{table}

These flags allow the policy to distinguish legal lane changes (dashed marking) from illegal
solid-line crossings without relying on CARLA's lane-invasion event alone.

\section{Vehicle state (dimensions 732\textendash 733)}\label{app:obs-vehicle}

\begin{table}[H]
  \centering
  \caption{Vehicle state features.}
  \label{tab:obs-vehicle}
  \small
  \begin{tabular}{@{}cllp{6cm}@{}}
    \toprule
    \textbf{Idx.} & \textbf{Name} & \textbf{Range} & \textbf{Meaning} \\
    \midrule
    732 & \texttt{speed\_norm} & \([0, 1]\) & Vehicle speed normalised by a dynamic reference:
        \(\max(\text{speed\_limit} \times 1.5,\; 10\,\mathrm{km/h})\). \\
    733 & \texttt{steering}    & \([-1, 1]\) & Last steering command actually applied to the
        vehicle (after shield correction). \\
    \bottomrule
  \end{tabular}
\end{table}

\section{Route information (dimensions 734\textendash 738)}\label{app:obs-route}

The final \(5\) dimensions provide look-ahead information along the route graph.

\begin{table}[H]
  \centering
  \caption{Route features.}
  \label{tab:obs-route}
  \small
  \begin{tabular}{@{}cllp{6cm}@{}}
    \toprule
    \textbf{Idx.} & \textbf{Name} & \textbf{Range} & \textbf{Meaning} \\
    \midrule
    734 & \texttt{next\_wp\_angle\_norm} & \([-1, 1]\) & Signed angle to the next waypoint
        \(5\,\mathrm{m}\) ahead, normalised by \(180^\circ\). \\
    735 & \texttt{wp\_angle\_20m\_norm} & \([-1, 1]\) & Signed angle to the waypoint
        \(20\,\mathrm{m}\) ahead, normalised by \(180^\circ\) (look-ahead for smooth turning). \\
    736 & \texttt{progress\_norm}       & \([0, 1]\) & Fraction of the success distance
        traversed so far during the current episode. \\
    737 & \texttt{speed\_limit\_norm}   & \([0, 1]\) & Current speed limit normalised by the
        reference speed. \\
    738 & \texttt{speed\_ratio}         & \([0, 2]\) & Ratio of the vehicle speed to the
        current speed limit (clipped at \(2\)). \\
    \bottomrule
  \end{tabular}
\end{table}

\section{Summary}\label{app:obs-summary}

The full observation vector has the structure
\begin{equation*}
  s_t \;=\; \underbrace{\big[\, \text{LIDAR}^{\text{comb}}_t \,\big]}_{240}
         \;\|\; \underbrace{\big[\, \text{LIDAR}^{\text{dyn}}_t \,\big]}_{240}
         \;\|\; \underbrace{\big[\, \text{LIDAR}^{\text{stat}}_t \,\big]}_{240}
         \;\|\; \underbrace{\big[\, \text{Lane}_t \,\big]}_{8}
         \;\|\; \underbrace{\big[\, \text{Marking}_t \,\big]}_{4}
         \;\|\; \underbrace{\big[\, \text{Veh}_t \,\big]}_{2}
         \;\|\; \underbrace{\big[\, \text{Route}_t \,\big]}_{5}
         \;\in\; \mathbb{R}^{739}.
\end{equation*}
Only the rightmost \(19\) dimensions (lane, marking, vehicle and route) are normalised online
by the running-mean estimator of Section~\ref{sec:design-obsnorm}; the \(720\) LIDAR dimensions
are already in \([0,1]\) by construction and are left untouched.
